% TOP_PROBLEM: {"id": 30006873, "title": "Log-Brunn-Minkowski Inequality for Symmetric Convex Bodies", "category_id": 6, "category": {"id": 6, "name": "geometry", "display_name": "Geometry", "description": "Euclidean and non-Euclidean geometry, geometric structures.", "slug": "geometry", "order_index": 6, "created_at": "2026-07-31T15:26:25.671Z"}, "status": "open"}
% FIRST_POSTED: 2026-09-27
% !TeX program = lualatex
% ATTEMPT_STATUS: unresolved
\documentclass[11pt]{article}
\usepackage[margin=25mm]{geometry}
\usepackage{fontspec}
\setmainfont{Latin Modern Roman}
\usepackage{amsmath,amssymb,mathtools}
\usepackage{unicode-math}
\setmathfont{Latin Modern Math}
\usepackage{xurl,hyperref}
\hypersetup{colorlinks=true,urlcolor=blue}
\setcounter{secnumdepth}{0}
\setlength{\parindent}{0pt}
\setlength{\parskip}{5pt}
\setlength{\emergencystretch}{3em}
\begin{document}
\section*{\texorpdfstring{Log-Brunn-Minkowski Inequality for Symmetric Convex Bodies}{Log-Brunn-Minkowski Inequality for Symmetric Convex Bodies}}\label{30006873}
\subsection{Definitions and mathematical statement}
For an origin-symmetric convex body $K\subset{\mathbb{R}}^n$, let $h_K(u)=\max_{x\in K}\langle x,u\rangle$ be its positive support function on $S^{n-1}$. Define its logarithmic combination with another such body $L$ by
\[
 (1-\lambda)K+_0\lambda L
 =\{x:\langle x,u\rangle\le
 h_K(u)^{1-\lambda}h_L(u)^\lambda\ \text{for all }u\in S^{n-1}\}.
\]
The log-Brunn--Minkowski conjecture asserts, for every $0\le\lambda\le1$,
\[
 |(1-\lambda)K+_0\lambda L|\ge |K|^{1-\lambda}|L|^\lambda.
\]
The intersection-of-halfspaces definition matters because the product of support functions need not itself be a support function. This is not the ordinary Minkowski sum $(1-\lambda)K+\lambda L$; symmetry about the origin is part of the hypotheses.
\par\smallskip\textit{Formulation and concept references:} \href{https://arxiv.org/abs/1311.4954}{[S1]}.

\subsection{Short English statement}
Does the logarithmic combination of two origin-symmetric convex bodies have volume at least the corresponding geometric mean of their volumes?
\subsection{Research attempt: exact boxes, unconditional bodies, and a local variation}
\paragraph{1. Scope and current-source audit.}
Write $W_\lambda(K,L)$ for the intersection of halfspaces in the question. The defining positive function $h_K^{1-\lambda}h_L^\lambda$ need not itself be a support function. This distinction will be checked explicitly below. Saroglou's paper [S1] treats unconditional bodies and related consequences. A current primary-source check also finds Lu's August 2026 preprint [E1], whose Theorems 1.2--1.3 assume $n-2$ common orthogonal reflection symmetries. That assumption is stronger than central symmetry in general; the theorem is not a solution for arbitrary pairs. Its introduction also states the known planar result. The independent calculations here cover boxes, recover the unconditional argument using the standard Pr\'ekopa--Leindler inequality, and establish an infinitesimal inequality at a disk. They make no novelty claim.

\paragraph{2. Affine covariance and endpoints.}
For an invertible linear map $A$, $h_{AK}(u)=h_K(A^Tu)$. The halfspace definition, using all nonzero directions and positive homogeneity, therefore gives
\[
 W_\lambda(AK,AL)=A W_\lambda(K,L).
\]
Both sides of the desired volume inequality multiply by $|\det A|$. A common linear change of coordinates is legitimate; applying unrelated changes to $K$ and $L$ is not. For $\lambda=0$ or $1$, the intersection of all supporting halfspaces returns $K$ or $L$, so equality holds. If $L=cK$, the defining function is $c^\lambda h_K$, and $W_\lambda=c^\lambda K$, again giving equality.

\paragraph{3. A complete calculation for boxes and the support-function trap.}
Let
\[
 K=\prod_{j=1}^n[-a_j,a_j],\qquad
 L=\prod_{j=1}^n[-b_j,b_j],\qquad a_j,b_j>0.
\]
The constraints in the coordinate directions show that $W_\lambda$ is contained in the box with half-side lengths $c_j=a_j^{1-\lambda}b_j^\lambda$. Conversely, if $|x_j|\leq c_j$, H\"older's inequality gives, for every $u\in\mathbb R^n$,
\[
 \langle x,u\rangle\leq\sum_j |u_j|a_j^{1-\lambda}b_j^\lambda
 \leq\left(\sum_j|u_j|a_j\right)^{1-\lambda}
       \left(\sum_j|u_j|b_j\right)^\lambda.
\]
The last expression is the defining bound. Hence the containment is equality, and
\[
 |W_\lambda|=2^n\prod_j a_j^{1-\lambda}b_j^\lambda
 =|K|^{1-\lambda}|L|^\lambda.
\]
This includes parallelepipeds sharing the same edge directions after a common linear change of coordinates.

Now take $n=2$, $a>1$, $K=[-a,a]\times[-1,1]$, $L=[-1,1]\times[-a,a]$, and $\lambda=1/2$. The homogeneous defining function $f=(h_Kh_L)^{1/2}$ satisfies
\[
 f(e_1)=f(e_2)=\sqrt a,\qquad f(e_1+e_2)=a+1>2\sqrt a.
\]
It violates subadditivity and thus cannot be a support function. Nevertheless the Wulff intersection is exactly $[-\sqrt a,\sqrt a]^2$. Its support at $e_1+e_2$ is $2\sqrt a$, strictly smaller than $f(e_1+e_2)$. Any proof differentiating $f$ as though it were the support of the intersection, with no convexity check, fails already on these elementary bodies.

\paragraph{4. The unconditional case with the key inclusion proved.}
Suppose $K$ and $L$ are invariant under changing the sign of any coordinate. If $x\in K$ has nonnegative coordinates, every vector $z$ satisfying $0\leq z_j\leq x_j$ also belongs to $K$. To see this, average a point with its reflection in one coordinate hyperplane to reduce that coordinate, and iterate. This downward-closed property uses both convexity and the individual reflections; central symmetry alone does not provide it.

For $x\in K\cap(0,\infty)^n$ and $y\in L\cap(0,\infty)^n$, put $z_j=x_j^{1-\lambda}y_j^\lambda$. For any $u$, unconditionality and H\"older give
\[
 \langle z,u\rangle\leq\sum_j|u_j|z_j
 \leq\langle x,|u|\rangle^{1-\lambda}
       \langle y,|u|\rangle^\lambda
 \leq h_K(u)^{1-\lambda}h_L(u)^\lambda.
\]
Thus $z\in W_\lambda$. Pass to logarithmic coordinates and define
\[
 A=\{s:e^s\in K\cap(0,\infty)^n\},\quad
 B=\{t:e^t\in L\cap(0,\infty)^n\},\quad
 D=\{v:e^v\in W_\lambda\cap(0,\infty)^n\},
\]
where exponentiation is coordinatewise. The inclusion just proved says $(1-\lambda)A+\lambda B\subset D$. Apply the standard Pr\'ekopa--Leindler integral inequality to
\[
 f(s)=\mathbf1_A(s)e^{\sum s_j},\quad
 g(t)=\mathbf1_B(t)e^{\sum t_j},\quad
 q(v)=\mathbf1_D(v)e^{\sum v_j}.
\]
For every $s,t$, $q((1-\lambda)s+\lambda t)\geq f(s)^{1-\lambda}g(t)^\lambda$; when either factor is zero the inequality is automatic. The change of variables $x_j=e^{s_j}$ identifies their integrals with the volumes in the positive orthant. Pr\'ekopa--Leindler therefore yields the desired geometric-mean inequality for these orthant volumes. The three bodies are unconditional, so each has $2^n$ times its positive-orthant volume; the factors cancel correctly because $(2^n)^{1-\lambda}(2^n)^\lambda=2^n$.

This proves the entire unconditional subcase, conditional only on the named, established integral inequality. It does not require falsely assuming that the defining geometric mean is a support function. For merely centrally symmetric bodies, the replacement of $u$ by $|u|$ and the independent reflection argument are unavailable. That is the exact point at which this proof stops.

\paragraph{5. A product inclusion and what it does not imply.}
Let $K=K_1\times K_2$ and $L=L_1\times L_2$ in a fixed direct-sum decomposition. If $x_i\in W_\lambda(K_i,L_i)$, then
\[
 \langle(x_1,x_2),(u_1,u_2)\rangle
 \leq\sum_{i=1}^2 h_{K_i}(u_i)^{1-\lambda}h_{L_i}(u_i)^\lambda
 \leq h_K(u_1,u_2)^{1-\lambda}h_L(u_1,u_2)^\lambda.
\]
Hence the product of the two lower-dimensional logarithmic combinations lies in the full combination. If the conjecture is known for both pairs of factors, multiplication of their volume inequalities proves it for the product pair. The result extends by induction to any fixed number of factors. General convex bodies do not decompose as Cartesian products, so this is a closure rule rather than a dimension induction for all bodies.

\paragraph{6. Second variation at a disk.}
For a smooth strictly convex planar body with support function $h(\theta)$ and positive curvature radius $h+h''$, its area is
\[
 A(h)=\frac12\int_0^{2\pi}(h^2-(h')^2)\,d\theta.
\]
Consider the logarithmic variation $h_t=h e^{tf}$, assumed to stay strictly convex for sufficiently small $|t|$. At $t=0$, differentiation and integration by parts give
\[
 A'=\int_0^{2\pi} f h(h+h'')\,d\theta,
 \qquad
 A''=\int_0^{2\pi}\bigl(2f^2h(h+h'')-h^2(f')^2\bigr)\,d\theta.
\]
One can check the second identity by differentiating the quadratic area functional twice: the $f^2$ term from $h_t''$ and the derivative of $fh$ combine, and the mixed term integrates to the displayed expression. Logarithmic concavity at $t=0$ means $AA''-(A')^2\leq0$.

For a disk, $h=R$ is constant. Central symmetry requires $f(\theta+\pi)=f(\theta)$. Write $f=c+g$, where $\int_0^{2\pi}g=0$. Then
\[
 AA''-(A')^2
 =\pi R^4\left(2\int_0^{2\pi}g^2-\int_0^{2\pi}(g')^2\right).
\]
The Fourier series of the $\pi$-periodic mean-zero function $g$ contains only frequencies $2,4,6,\ldots$, so Parseval gives $\int(g')^2\geq4\int g^2$. Thus the second variation is nonpositive, and strictly negative unless $f$ is constant. Constant $f$ describes dilation and gives equality, as it should.

This calculation proves a local infinitesimal statement at the disk. It does not prove the global inequality by integrating along $h_K^{1-t}h_L^t$, because that path may fail to consist of support functions, and the calculation at one body is not a bound at every body along a path. Both failures are independent and must be addressed in any global argument.

\paragraph{7. Exact diagnostics and outcome.}
The saved diagnostic checks the box support-function failure at rational $a=4$, the exact box volume identity for selected rational side lengths at $\lambda=1/2$ with square products, and the Fourier coefficients $2-(2m)^2$ in the disk second variation. These are finite algebraic checks. They do not test all centrally symmetric convex bodies, all support directions, or a continuum of logarithmic paths.

\textbf{UNRESOLVED.} The unconditional theorem, exact box equality, product closure, and disk second variation are established with their stated inputs. The general case remains outside these arguments because independent coordinate sign symmetries are absent and the Wulff envelope can differ from the proposed support interpolation.

\subsection{Sources}
\begin{itemize}
\item[S1]C. Saroglou, Remarks on the conjectured log-Brunn-Minkowski inequality, Geometriae Dedicata 172 (2014) 405-418; arXiv:1311.4954.
\url{https://arxiv.org/abs/1311.4954}.
\textit{Formulation source.}
\item[E1]X. Lu, Logarithmic Brunn--Minkowski Inequality under $n-2$ Reflection Symmetries, arXiv:2608.20730v2, 27 August 2026, Theorems 1.2--1.3.
\url{https://arxiv.org/html/2608.20730v2}.
\textit{Scope checked 27 September 2026: common orthogonal reflection assumptions, not arbitrary centrally symmetric pairs. The new paper's full proof is not independently audited here.}
\end{itemize}
\end{document}
